\documentclass[10pt,letterpaper]{article}
\usepackage[T1]{fontenc}
\usepackage[utf8]{inputenc}
\usepackage{lmodern}
\usepackage[margin=0.83in,headheight=16pt]{geometry}
\usepackage{array,tabularx,longtable,booktabs,xcolor}
\usepackage{enumitem,fancyhdr,hyperref,lastpage}
\definecolor{navy}{HTML}{19334D}
\definecolor{teal}{HTML}{146B74}
\definecolor{muted}{HTML}{596B7B}
\definecolor{mist}{HTML}{EAF1F3}
\hypersetup{hidelinks,pdfauthor={Enterprise Architecture Documentation},pdfsubject={Enterprise technology architecture and operations}}
\setlength{\parindent}{0pt}
\setlength{\parskip}{6pt}
\setlist[itemize]{topsep=3pt,itemsep=3pt,leftmargin=16pt}
\setlist[enumerate]{topsep=3pt,itemsep=3pt,leftmargin=19pt}
\renewcommand{\arraystretch}{1.23}
\newcolumntype{Y}{>{\raggedright\arraybackslash}X}
\pagestyle{fancy}\fancyhf{}
\fancyhead[L]{\scriptsize\color{teal} ENTERPRISE TECHNOLOGY | ARCHITECTURE OPERATIONS}
\fancyhead[R]{\scriptsize\color{muted} 2026.10}
\fancyfoot[L]{\scriptsize\color{muted} Proposed reference; not production approval}
\fancyfoot[R]{\scriptsize\thepage\ / \pageref{LastPage}}
\renewcommand{\headrulewidth}{0.4pt}
\newcommand{\writearea}[1][0.53in]{\noindent\fbox{\parbox{\dimexpr\textwidth-2\fboxsep-2\fboxrule\relax}{\rule{0pt}{#1}\hspace{0.95\linewidth}}}\par\vspace{0.22cm}}
\begin{document}
\begin{center}
{\large\bfseries\color{teal} ENTERPRISE ARCHITECTURE DOCUMENTATION SUITE}\par
\vspace{9pt}{\LARGE\bfseries\color{navy} Capability Overlap, Options and Tool Selection}\par
\vspace{8pt}{\small\color{muted} EOT-REF-02 | Revision 1.0 | 31 October 2026}\par
\end{center}
\vspace{8pt}\hrule\vspace{12pt}
\textbf{Purpose and status.} Decide what is complementary, competing or outside scope before accepting operational overhead.\par
\textbf{Evidence status:} Proposed / requires local validation.\par
\section{Overlap and complement decision register}
\small\begin{longtable}{>{\raggedright\arraybackslash}p{0.19\linewidth}>{\raggedright\arraybackslash}p{0.31\linewidth}>{\raggedright\arraybackslash}p{0.44\linewidth}}
\toprule
\textbf{Capabilities} & \textbf{Overlap or relationship} & \textbf{Disposition}\\
\midrule\endfirsthead
\toprule
\textbf{Capabilities} & \textbf{Overlap or relationship} & \textbf{Disposition}\\
\midrule\endhead
GLPI / iTop & Both can manage asset/service data and tickets & Assign GLPI inventory; iTop service/CI/ITSM. Or select one to reduce cost.\\
GHAS / Semgrep OSS & Overlapping SAST; GHAS is not OSS & Default strict-OSS option Semgrep OSS; commercial exception review for GHAS.\\
Trivy / Grype & Image and package vulnerability overlap & Select a primary scanner; retain secondary only to test coverage.\\
Syft / Trivy / Dependency-Track & SBOM generation plus analysis & Syft output and Dependency-Track lifecycle complement scanning; avoid duplicate billing.\\
DefectDojo / Taiga / iTop & Case triage, stories, incidents can overlap & Single findings master; Taiga delivery; iTop incidents/changes.\\
Authentik / Keycloak & SSO and federation overlap & Choose one IdP after SSO/HA/maintenance analysis.\\
MISP / OpenCTI & Threat intelligence platform overlap & Not baseline; deploy to proven use cases.\\
GnuCash / Actual / Firefly III & Ledger vs budgeting/personal reporting & No substitution without accounting requirements assessment.\\
\bottomrule
\end{longtable}
\normalsize
\section{Selection hard gates}
\begin{itemize}
\item License: inspect all deployed edition/components for OSI-approved or approved commercial exception status.
\item Data: classify inputs, residency, retention, deletion and export obligations before integration.
\item Security: maintenance activity, vulnerability response and secret storage must be evidenced.
\item Operations: upgrade, backup, restore, RTO/RPO and health check must be demonstrated.
\item Function: demonstrate required APIs/SSO feature in chosen release; do not infer from marketing language.
\end{itemize}
\section{Comparative scoring method}
Apply 1-5 scores with evidence for functionality 25\%, maintenance 20\%, security 15\%, integration 15\%, operating burden 10\%, license fit 10\%, documentation 5\%. Weighted scoring compares passed candidates only; unknown hard-gate evidence is HOLD rather than a numerical zero.\par
\textbf{Decision / gate.} Avoid claims that user preference, GitHub stars, or license status alone proves enterprise readiness.\par
\section{Candidate decision examples}
\small\begin{longtable}{>{\raggedright\arraybackslash}p{0.19\linewidth}>{\raggedright\arraybackslash}p{0.31\linewidth}>{\raggedright\arraybackslash}p{0.44\linewidth}}
\toprule
\textbf{Need} & \textbf{Proposed baseline} & \textbf{Conditional alternatives}\\
\midrule\endfirsthead
\toprule
\textbf{Need} & \textbf{Proposed baseline} & \textbf{Conditional alternatives}\\
\midrule\endhead
Secret detection & Gitleaks & GHAS secret protection under approved license\\
SAST & Semgrep OSS & GHAS code security or language-specific analyzers\\
Web DAST & ZAP in staging & Manual targeted testing\\
Package vulnerability & Trivy & Grype for second opinion\\
Finding triage & Existing iTop case or DefectDojo & Custom service only if justified\\
Secure backup & restic with restore evidence & Provider-native encrypted backups\\
Personal budgeting & Actual or HomeBank & Firefly III or MMEX\\
Small-business ledger & GnuCash after accountant review & Dedicated accounting product outside scope\\
\bottomrule
\end{longtable}
\normalsize
\vfill\noindent\color{muted}\rule{\linewidth}{0.4pt}\par\scriptsize Original synthesis of three repository sources; tool-specific integration and license claims require edition/version-specific verification.\normalsize
\end{document}
